\section{Open Challenges and Roadmap}
\label{sec:roadmap}

\subsection{From quantum-inspired to true quantum}
The quantum-inspired operator runs on classical hardware because it is a matrix
function of the Laplacian. Realizing the walk on quantum hardware promises the
exponential-dimensional feature maps a classical model cannot represent, but several
obstacles stand in the way. Noisy intermediate-scale quantum (NISQ) devices~\cite{preskill2018nisq} introduce
depolarizing and readout noise that degrades accuracy; qubit counts cap the graph size
that can be embedded directly; and parameterized quantum
circuits~\cite{cerezo2021vqa} suffer barren
plateaus~\cite{mcclean2018barren} where gradients vanish as the system grows. Quantifying how much of the
ballistic mechanism survives device noise is a concrete, near-term experiment: a small
parameterized quantum walk on a noisy simulator, swept over depolarizing strength, to
measure where the far-node edge of Section~\ref{sec:casestudy} erodes. We leave this to
future work.

\subsection{Scaling the classical surrogate}
Even as a classical surrogate, the eigendecomposition behind $e^{-it\Lap}$ costs
$O(N^3)$, which is acceptable at the hundreds-of-nodes scale of this tutorial but not
for full mega-constellations. Krylov-subspace and Chebyshev-polynomial approximations
compute $e^{-it\Lap}\mathbf{x}$ without a full spectrum and are the practical route to
thousands of nodes; characterizing the accuracy-cost tradeoff of these approximations
for quantum-walk kernels is an open engineering problem.

\subsection{Beyond a single snapshot}
LEO topologies are not only dynamic but \emph{predictably} dynamic: the future graph
is known from ephemeris. The dynamic-graph methods surveyed in
Section~\ref{sec:survey}, TGN~\cite{rossi2020tgn}, EvolveGCN~\cite{pareja2020evolvegcn},
and DySAT~\cite{sankar2020dysat}, learn over a sequence of \emph{past} snapshots and do
not use a known future; an operator that propagates over a short horizon of future
Laplacians rather than a single snapshot is a promising direction that the snapshot-based
study here does not cover.

\subsection{Federated and non-centralized quantum graph learning}
In a constellation the graph is physically distributed: each satellite or region holds
part of the topology and its local traffic. On-board and cross-satellite federated
learning is already an active line for LEO
networks~\cite{razmi2022board,matthiesen2024federated}, building on federated
optimization~\cite{mcmahan2017federated}. Running quantum or quantum-inspired graph
operators in this setting, where no node sees the whole Laplacian, is open: the spectral
operators of this tutorial assume a global eigenbasis, so decentralized or
locally-approximated propagation is required. The non-centralized processing of both
classical data and quantum states, an explicit theme of this special issue, meets the
quantum-internet roadmap here~\cite{wehner2018quantum}.

\subsection{Quantum algorithms for network functions}
Beyond learning, quantum search and optimization target network functions directly, for
example virtualization and resource placement~\cite{botsinis2019quantum}. How these
combine with the learned propagation operators studied here, a learned model proposing
candidates that a quantum subroutine refines, is an unexplored but natural hybrid for
future wireless management.

\subsection{Benchmarks and reproducibility}
Finally, the field needs shared, scale-free benchmarks for graph learning on dynamic
wireless topologies, with multi-seed protocols and packet-level
validation~\cite{kassing2020hypatia} built in. The pitfalls of
Section~\ref{sec:pitfalls} persist largely because evaluation is ad hoc, an observation
echoed across machine
learning~\cite{henderson2018matters,pineau2021reproducibility,dehghani2021lottery}. The
companion code to this tutorial is a small step in that direction: every figure is
regenerated by a single script from seeded runs.
